
\documentclass{article}
\usepackage{colortbl}
\usepackage{makecell}
\usepackage{multirow}
\usepackage{supertabular}

\begin{document}

\newcounter{utterance}

\centering \large Interaction Transcript for game `matchit\_ascii', experiment `same\_grid', episode 3 with reflexion{-}llama{-}3.1{-}8b{-}instruct.
\vspace{24pt}

{ \footnotesize  \setcounter{utterance}{1}
\setlength{\tabcolsep}{0pt}
\begin{supertabular}{c@{$\;$}|p{.15\linewidth}@{}p{.15\linewidth}p{.15\linewidth}p{.15\linewidth}p{.15\linewidth}p{.15\linewidth}}
    \# & $\;$A & \multicolumn{4}{c}{Game Master} & $\;\:$B\\
    \hline

    \theutterance \stepcounter{utterance}  
    & & \multicolumn{4}{p{0.6\linewidth}}{
        \cellcolor[rgb]{0.9,0.9,0.9}{
            \makecell[{{p{\linewidth}}}]{
                \texttt{\tiny{[P1$\langle$GM]}}
                \texttt{You have attempted this same task before and did not succeed. These are the plans you wrote after those attempts.} \\
\texttt{Trial \#1: New plan: } \\
\\ 
\texttt{Upon reflection, I realize that I didn't ask enough targeted questions to gather crucial information about the other grid. I got stuck in a loop of asking yes/no questions that didn't provide much insight into the grid's structure. } \\
\\ 
\texttt{Next time, I will ask more open{-}ended questions that can help me understand the grid's layout and relationships between symbols. For example, I could ask about the number of X symbols in a specific row or column, or the distance between two adjacent X symbols. I will also pay closer attention to the answers provided and use that information to inform my next questions.} \\
\\ 
\texttt{Trial \#2: New plan: } \\
\\ 
\texttt{Upon reflection, I realize that I didn't effectively use the information gathered from the other player's answers to inform my next questions. I got stuck in a loop of asking similar questions without making progress towards understanding the grid's structure.} \\
\\ 
\texttt{Next time, I will pay closer attention to the answers provided and use that information to ask more targeted and insightful questions. For example, when the other player mentioned that the third row in their grid is completely filled with ▢ symbols, I should have asked follow{-}up questions to understand the implications of this information, such as whether this row is also present in my grid or if there are any other rows with a similar property. I will also try to ask questions that can help me understand the relationships between symbols, such as the number of X symbols adjacent to a specific ▢ symbol or the number of ▢ symbols between two adjacent X symbols.} \\
\\ 
\texttt{Trial \#3: New plan: } \\
\\ 
\texttt{Upon reflection, I realize that I didn't effectively use the information gathered from the other player's answers to inform my next questions, and I also didn't ask enough targeted questions to gather crucial information about the other grid. I got stuck in a loop of asking similar questions without making progress towards understanding the grid's structure.} \\
\\ 
\texttt{Next time, I will pay closer attention to the answers provided and use that information to ask more targeted and insightful questions. I will also ask more open{-}ended questions that can help me understand the grid's layout and relationships between symbols. } \\
\\ 
\texttt{To achieve this, I will follow a more structured approach:} \\
\\ 
\texttt{1.  After receiving an answer, I will take a moment to analyze the information and think about the implications.} \\
\texttt{2.  I will then ask follow{-}up questions to clarify or gather more information about the answer.} \\
\texttt{3.  I will also try to ask questions that can help me understand the relationships between symbols, such as the number of X symbols adjacent to a specific ▢ symbol or the number of ▢ symbols between two adjacent X symbols.} \\
\texttt{4.  I will also pay attention to the answers provided and use that information to inform my next questions, rather than asking similar questions without making progress.} \\
\\ 
\texttt{By following this approach, I hope to gather more crucial information about the other grid and make more informed decisions about the grid's structure.} \\
\\ 
\\ 
\texttt{You are participating in a collaborative guessing game.} \\
\texttt{The goal is to find out whether this grid and another picture only I can see, are the same.} \\
\texttt{Please describe your grid first. Then, I will provide my description and we can ask each other questions about the grids to figure out whether they are the same. Now start your short grid description with "DESCRIPTION:" followed by the description. Do not add anything else. This is your grid:} \\
\texttt{X ▢ ▢ ▢ X} \\
\texttt{▢ X ▢ X ▢} \\
\texttt{▢ ▢ X ▢ ▢} \\
\texttt{▢ X ▢ X ▢} \\
\texttt{X ▢ ▢ ▢ X} \\
            }
        }
    }
    & \\ \\

    \theutterance \stepcounter{utterance}  
    & \multicolumn{4}{p{0.6\linewidth}}{
        \cellcolor[rgb]{0.9,0.9,0.9}{
            \makecell[{{p{\linewidth}}}]{
                \texttt{\tiny{[P1$\rangle$GM]}}
                \texttt{DESCRIPTION: The grid is a 5x5 grid with a mix of X and ▢ symbols. The X symbols are located at the corners and the center of the grid, with two X symbols in the first and last rows. The ▢ symbols are scattered throughout the grid, with some rows and columns containing more ▢ symbols than others.} \\
            }
        }
    }
    & & \\ \\

    \theutterance \stepcounter{utterance}  
    & & & \multicolumn{2}{p{0.3\linewidth}}{
        \cellcolor[rgb]{0.9,0.9,0.9}{
            \makecell[{{p{\linewidth}}}]{
                \texttt{\tiny{[GM$|$GM]}}
                \texttt{continue} \\
            }
        }
    }
    & & \\ \\

    \theutterance \stepcounter{utterance}  
    & & \multicolumn{4}{p{0.6\linewidth}}{
        \cellcolor[rgb]{0.9,0.9,0.9}{
            \makecell[{{p{\linewidth}}}]{
                \texttt{\tiny{[GM$\rangle$P2]}}
                \texttt{You have attempted this same task before and did not succeed. These are the plans you wrote after those attempts.} \\
\texttt{Trial \#1: New plan: } \\
\\ 
\texttt{I got stuck in a loop of asking questions that didn't provide enough new information to make a decision. I started with a good question about the diagonal pattern, but then I asked several questions that were too specific and didn't reveal any significant differences between the grids. } \\
\\ 
\texttt{Next time, I will start with a more general question that can help me understand the overall structure of the grid, and then ask more specific questions based on the answers I receive. I will also try to ask questions that can help me identify any unique features of the grid that might indicate whether it's the same or different from my grid. } \\
\\ 
\texttt{For example, I could ask a question like "Are there any rows or columns in your grid that have a specific number of X symbols?" or "Are there any X symbols in your grid that are not part of the diagonal pattern?" These types of questions can help me gather more information and make a more informed decision about whether the grids are the same or different.} \\
\\ 
\texttt{Trial \#2: New plan: } \\
\\ 
\texttt{I got stuck in a loop of asking questions that didn't provide enough new information to make a decision. I started with a good question about the diagonal pattern, but then I asked several questions that were too specific and didn't reveal any significant differences between the grids. } \\
\\ 
\texttt{Next time, I will start with a more general question that can help me understand the overall structure of the grid, and then ask more specific questions based on the answers I receive. I will also try to ask questions that can help me identify any unique features of the grid that might indicate whether it's the same or different from my grid. } \\
\\ 
\texttt{However, I also failed to make a decision at the end, despite having some information about the differences between the grids. Next time, I will make sure to carefully consider the answers to my questions and make a decision based on the information I have gathered. } \\
\\ 
\texttt{Specifically, I will pay closer attention to the answers about the third row, which seemed to be a key difference between the grids. I will also try to ask questions that can help me understand the relationship between the X symbols and the ▢ symbols in the grid, as this seemed to be an important aspect of the game. } \\
\\ 
\texttt{For example, I could ask a question like "Are there any ▢ symbols in your grid that are adjacent to two X symbols?" or "Are there any rows or columns in your grid that have a specific number of ▢ symbols?" These types of questions can help me gather more information and make a more informed decision about whether the grids are the same or different.} \\
\\ 
\texttt{Trial \#3: New plan: } \\
\\ 
\texttt{I got stuck in a loop of asking questions that didn't provide enough new information to make a decision. I started with a good question about the diagonal pattern, but then I asked several questions that were too specific and didn't reveal any significant differences between the grids. } \\
\\ 
\texttt{However, I also failed to make a decision at the end, despite having some information about the differences between the grids. Next time, I will make sure to carefully consider the answers to my questions and make a decision based on the information I have gathered. } \\
\\ 
\texttt{Specifically, I will pay closer attention to the answers about the third row, which seemed to be a key difference between the grids. I will also try to ask questions that can help me understand the relationship between the X symbols and the ▢ symbols in the grid, as this seemed to be an important aspect of the game. } \\
\\ 
\texttt{In addition, I will make sure to ask questions that can help me identify any unique features of the grid that might indicate whether it's the same or different from my grid. I will also try to ask more general questions that can help me understand the overall structure of the grid, and then ask more specific questions based on the answers I receive. } \\
\\ 
\texttt{For example, I could ask a question like "Are there any rows or columns in your grid that have a specific number of X symbols?" or "Are there any X symbols in your grid that are not part of the diagonal pattern?" These types of questions can help me gather more information and make a more informed decision about whether the grids are the same or different. } \\
\\ 
\texttt{I will also try to ask questions that can help me understand the distribution of X symbols and ▢ symbols in the grid, such as "Are there any rows or columns in your grid that have a specific number of ▢ symbols?" or "Are there any ▢ symbols in your grid that are adjacent to three X symbols?" These types of questions can help me identify any unique features of the grid that might indicate whether it's the same or different from my grid. } \\
\\ 
\texttt{Ultimately, my goal is to gather enough information to make a confident decision about whether the grids are the same or different, and I will continue to ask questions and refine my strategy until I achieve this goal.} \\
\\ 
\\ 
\texttt{You are participating in a collaborative guessing game.} \\
\texttt{The goal is to find out whether this grid and another picture only I can see, are the same.} \\
\texttt{Please describe your grid first. Then, I will provide my description and we can ask each other questions about the grids to figure out whether they are the same. Now start your short grid description with "DESCRIPTION:" followed by the description. Do not add anything else. This is your grid:} \\
\texttt{X ▢ ▢ ▢ X} \\
\texttt{▢ X ▢ X ▢} \\
\texttt{▢ ▢ X ▢ ▢} \\
\texttt{▢ X ▢ X ▢} \\
\texttt{X ▢ ▢ ▢ X} \\
            }
        }
    }
    & \\ \\

    \theutterance \stepcounter{utterance}  
    & & & \multicolumn{4}{p{0.6\linewidth}}{
        \cellcolor[rgb]{0.9,0.9,0.9}{
            \makecell[{{p{\linewidth}}}]{
                \texttt{\tiny{[GM$\langle$P2]}}
                \texttt{DESCRIPTION: The grid is a 5x5 grid with a diagonal pattern of X symbols. The X symbols are located at the top{-}left and bottom{-}right corners, and the middle of the grid. The ▢ symbols are scattered throughout the grid, with some rows and columns having a specific number of ▢ symbols.} \\
            }
        }
    }
     \\ \\

    \theutterance \stepcounter{utterance}  
    & & & \multicolumn{2}{p{0.3\linewidth}}{
        \cellcolor[rgb]{0.9,0.9,0.9}{
            \makecell[{{p{\linewidth}}}]{
                \texttt{\tiny{[GM$|$GM]}}
                \texttt{continue} \\
            }
        }
    }
    & & \\ \\

    \theutterance \stepcounter{utterance}  
    & & \multicolumn{4}{p{0.6\linewidth}}{
        \cellcolor[rgb]{0.9,0.9,0.9}{
            \makecell[{{p{\linewidth}}}]{
                \texttt{\tiny{[GM$\rangle$P2]}}
                \texttt{This is my DESCRIPTION: The grid is a 5x5 grid with a mix of X and ▢ symbols. The X symbols are located at the corners and the center of the grid, with two X symbols in the first and last rows. The ▢ symbols are scattered throughout the grid, with some rows and columns containing more ▢ symbols than others.} \\
\texttt{Now ask a question in order to find out new aspects of my grid that may be different to your grid. Start with "QUESTION:" and do not add anything else.} \\
            }
        }
    }
    & \\ \\

    \theutterance \stepcounter{utterance}  
    & & & \multicolumn{4}{p{0.6\linewidth}}{
        \cellcolor[rgb]{0.9,0.9,0.9}{
            \makecell[{{p{\linewidth}}}]{
                \texttt{\tiny{[GM$\langle$P2]}}
                \texttt{QUESTION: Are there any rows or columns in your grid that have a specific number of ▢ symbols, and if so, which ones?} \\
            }
        }
    }
     \\ \\

    \theutterance \stepcounter{utterance}  
    & & & \multicolumn{2}{p{0.3\linewidth}}{
        \cellcolor[rgb]{0.9,0.9,0.9}{
            \makecell[{{p{\linewidth}}}]{
                \texttt{\tiny{[GM$|$GM]}}
                \texttt{continue} \\
            }
        }
    }
    & & \\ \\

    \theutterance \stepcounter{utterance}  
    & & \multicolumn{4}{p{0.6\linewidth}}{
        \cellcolor[rgb]{0.9,0.9,0.9}{
            \makecell[{{p{\linewidth}}}]{
                \texttt{\tiny{[P1$\langle$GM]}}
                \texttt{This is my DESCRIPTION: The grid is a 5x5 grid with a diagonal pattern of X symbols. The X symbols are located at the top{-}left and bottom{-}right corners, and the middle of the grid. The ▢ symbols are scattered throughout the grid, with some rows and columns having a specific number of ▢ symbols.} \\
\texttt{QUESTION: Are there any rows or columns in your grid that have a specific number of ▢ symbols, and if so, which ones? Start your answer with "ANSWER:" and do not add anything else.} \\
            }
        }
    }
    & \\ \\

    \theutterance \stepcounter{utterance}  
    & \multicolumn{4}{p{0.6\linewidth}}{
        \cellcolor[rgb]{0.9,0.9,0.9}{
            \makecell[{{p{\linewidth}}}]{
                \texttt{\tiny{[P1$\rangle$GM]}}
                \texttt{ANSWER: Yes, the first and last rows in my grid have a specific number of ▢ symbols, with 3 ▢ symbols in each row.} \\
            }
        }
    }
    & & \\ \\

    \theutterance \stepcounter{utterance}  
    & & & \multicolumn{2}{p{0.3\linewidth}}{
        \cellcolor[rgb]{0.9,0.9,0.9}{
            \makecell[{{p{\linewidth}}}]{
                \texttt{\tiny{[GM$|$GM]}}
                \texttt{continue} \\
            }
        }
    }
    & & \\ \\

    \theutterance \stepcounter{utterance}  
    & & \multicolumn{4}{p{0.6\linewidth}}{
        \cellcolor[rgb]{0.9,0.9,0.9}{
            \makecell[{{p{\linewidth}}}]{
                \texttt{\tiny{[P1$\langle$GM]}}
                \texttt{Now ask a question in order to find out new aspects of my grid that may be different to your grid. Start with "QUESTION:" and do not add anything else.} \\
            }
        }
    }
    & \\ \\

    \theutterance \stepcounter{utterance}  
    & \multicolumn{4}{p{0.6\linewidth}}{
        \cellcolor[rgb]{0.9,0.9,0.9}{
            \makecell[{{p{\linewidth}}}]{
                \texttt{\tiny{[P1$\rangle$GM]}}
                \texttt{QUESTION: Are there any rows or columns in your grid that are completely filled with ▢ symbols?} \\
            }
        }
    }
    & & \\ \\

    \theutterance \stepcounter{utterance}  
    & & & \multicolumn{2}{p{0.3\linewidth}}{
        \cellcolor[rgb]{0.9,0.9,0.9}{
            \makecell[{{p{\linewidth}}}]{
                \texttt{\tiny{[GM$|$GM]}}
                \texttt{continue} \\
            }
        }
    }
    & & \\ \\

    \theutterance \stepcounter{utterance}  
    & & \multicolumn{4}{p{0.6\linewidth}}{
        \cellcolor[rgb]{0.9,0.9,0.9}{
            \makecell[{{p{\linewidth}}}]{
                \texttt{\tiny{[GM$\rangle$P2]}}
                \texttt{ANSWER: Yes, the first and last rows in my grid have a specific number of ▢ symbols, with 3 ▢ symbols in each row.} \\
\texttt{QUESTION: Are there any rows or columns in your grid that are completely filled with ▢ symbols? Start your answer with "ANSWER:" and do not add anything else.} \\
            }
        }
    }
    & \\ \\

    \theutterance \stepcounter{utterance}  
    & & & \multicolumn{4}{p{0.6\linewidth}}{
        \cellcolor[rgb]{0.9,0.9,0.9}{
            \makecell[{{p{\linewidth}}}]{
                \texttt{\tiny{[GM$\langle$P2]}}
                \texttt{QUESTION: Are there any rows or columns in your grid that are completely filled with X symbols?} \\
            }
        }
    }
     \\ \\

    \theutterance \stepcounter{utterance}  
    & & & \multicolumn{2}{p{0.3\linewidth}}{
        \cellcolor[rgb]{0.9,0.9,0.9}{
            \makecell[{{p{\linewidth}}}]{
                \texttt{\tiny{[GM$|$GM]}}
                \texttt{abort, first word: QUESTION:, but expected ANSWER:} \\
            }
        }
    }
    & & \\ \\

    \theutterance \stepcounter{utterance}  
    & \multicolumn{4}{p{0.6\linewidth}}{
        \cellcolor[rgb]{0.95,0.95,0.95}{
            \makecell[{{p{\linewidth}}}]{
                \texttt{\tiny{[P2$\rightarrow$P1]}}
                \texttt{Post{-}episode Reflexion (not sent to GM):} \\
\texttt{New plan: } \\
\\ 
\texttt{Upon reflection, I realize that I didn't effectively use the information gathered from the other player's answers to inform my next questions, and I also didn't ask enough targeted questions to gather crucial information about the other grid. I got stuck in a loop of asking similar questions without making progress towards understanding the grid's structure.} \\
\\ 
\texttt{Next time, I will pay closer attention to the answers provided and use that information to ask more targeted and insightful questions. I will also ask more open{-}ended questions that can help me understand the grid's layout and relationships between symbols. } \\
\\ 
\texttt{To achieve this, I will follow a more structured approach:} \\
\\ 
\texttt{1.  After receiving an answer, I will take a moment to analyze the information and think about the implications.} \\
\texttt{2.  I will then ask follow{-}up questions to clarify or gather more information about the answer.} \\
\texttt{3.  I will also try to ask questions that can help me understand the relationships between symbols, such as the number of X symbols adjacent to a specific ▢ symbol or the number of ▢ symbols between two adjacent X symbols.} \\
\texttt{4.  I will also pay attention to the answers provided and use that information to inform my next questions, rather than asking similar questions without making progress.} \\
\\ 
\texttt{However, I also realize that I didn't follow the format specified by the Game Master in my previous attempts. I will make sure to follow the format correctly in my next attempt, which means starting my answer with "ANSWER:" or "QUESTION:" as required.} \\
            }
        }
    }
    & & \\ \\

    \theutterance \stepcounter{utterance}  
    & & & \multicolumn{4}{p{0.6\linewidth}}{
        \cellcolor[rgb]{0.95,0.95,0.95}{
            \makecell[{{p{\linewidth}}}]{
                \texttt{\tiny{[P1$\rightarrow$P2]}}
                \texttt{Post{-}episode Reflexion (not sent to GM):} \\
\texttt{New plan: } \\
\\ 
\texttt{I got stuck in a loop of asking questions that didn't provide enough new information to make a decision. I started with a good question about the diagonal pattern, but then I asked several questions that were too specific and didn't reveal any significant differences between the grids. } \\
\\ 
\texttt{However, I also failed to make a decision at the end, despite having some information about the differences between the grids. Next time, I will make sure to carefully consider the answers to my questions and make a decision based on the information I have gathered. } \\
\\ 
\texttt{Specifically, I will pay closer attention to the answers about the third row, which seemed to be a key difference between the grids. I will also try to ask questions that can help me understand the relationship between the X symbols and the ▢ symbols in the grid, as this seemed to be an important aspect of the game. } \\
\\ 
\texttt{In addition, I will make sure to ask questions that can help me identify any unique features of the grid that might indicate whether it's the same or different from my grid. I will also try to ask more general questions that can help me understand the overall structure of the grid, and then ask more specific questions based on the answers I receive. } \\
\\ 
\texttt{However, I also failed to follow the format of the game, which led to the game being aborted. Next time, I will make sure to follow the format and ask questions in the correct format, starting with "QUESTION:" and not adding any other text. } \\
\\ 
\texttt{For example, I could ask a question like "QUESTION: Are there any rows or columns in your grid that have a specific number of ▢ symbols?" or "QUESTION: Are there any X symbols in your grid that are not part of the diagonal pattern?" These types of questions can help me gather more information and make a more informed decision about whether the grids are the same or different.} \\
            }
        }
    }
     \\ \\

\end{supertabular}
}

\end{document}
